% Cambio de base y homomorfismos

\begin{frame}
    \frametitle{La matriz de un homomorfismo depende de las bases}

    \begin{example}
        Sea $f\colon\R^2\to\R^2$, $f(u_1,u_2)=(2u_1+u_2,\ u_1+2u_2)$.
    \end{example}

    En la base canónica $C$, $f(e_1)=(2,1)$ y $f(e_2)=(1,2)$:
    \[
        M_{C,C}(f)=\begin{pmatrix}2&1\\1&2\end{pmatrix}.
    \]
    En la base $B=\{(1,1),(1,-1)\}$ ocurre algo distinto:
    \[
        f(1,1)=(3,3)=3\,(1,1),\qquad f(1,-1)=(1,-1)=1\,(1,-1),
    \]
    luego
    \[
        M_{B,B}(f)=\begin{pmatrix}3&0\\0&1\end{pmatrix}.
    \]

    \medskip

    Es \alert{el mismo} homomorfismo, representado por dos matrices
    distintas, y la segunda es mucho más sencilla. ¿Cómo se relacionan?
\end{frame}


\begin{frame}[allowframebreaks, t]
    \frametitle{Ejemplo: otra forma de calcular una matriz}

    \begin{example}
        Recalculemos $M_{B,B'}(f)$ para
        $f(u_1,u_2)=(u_1+u_2,\,2u_1-u_2,\,3u_2)$,
        $B=\{(1,1),(1,-1)\}$ y $B'=\{(1,0,0),(1,1,0),(1,1,1)\}$.
    \end{example}

    Con las bases canónicas $C$ de $\R^2$ y $C'$ de $\R^3$,
    $M_{C,C'}(f)=\begin{pmatrix}1&1\\2&-1\\0&3\end{pmatrix}$. Por la
    fórmula del cambio de base,
    \[
        M_{B,B'}(f)=M(C',B')\,M_{C,C'}(f)\,M(B,C),
    \]
    donde $M(B,C)=\begin{pmatrix}1&1\\1&-1\end{pmatrix}$ y
    $M(C',B')=M(B',C')^{-1}$, con
    \[
        M(B',C')=\begin{pmatrix}1&1&1\\0&1&1\\0&0&1\end{pmatrix},
        \qquad
        M(C',B')=\begin{pmatrix}1&-1&0\\0&1&-1\\0&0&1\end{pmatrix}.
    \]

    \framebreak

    Multiplicando,
    \begin{align*}
        M_{B,B'}(f)&=
        \begin{pmatrix}1&-1&0\\0&1&-1\\0&0&1\end{pmatrix}
        \begin{pmatrix}1&1\\2&-1\\0&3\end{pmatrix}
        \begin{pmatrix}1&1\\1&-1\end{pmatrix}\\
        &=
        \begin{pmatrix}1&-1&0\\0&1&-1\\0&0&1\end{pmatrix}
        \begin{pmatrix}2&0\\1&3\\3&-3\end{pmatrix}
        =
        \begin{pmatrix}1&-3\\-2&6\\3&-3\end{pmatrix},
    \end{align*}
    que coincide con el resultado obtenido antes resolviendo los sistemas.
\end{frame}


\begin{frame}
    \frametitle{Toda matriz invertible es una matriz de cambio de base}

    \begin{lemma}
        Sea $B$ una base de un espacio $V$ de dimensión $n$ y sea
        $P\in M_n(k)$ invertible. Existe una base $C$ de $V$ tal que
        $M(C,B)=P$.
    \end{lemma}

    \begin{proof}
        Como $c_B$ es sobreyectiva, existen $v'_1,\ldots,v'_n\in V$ tales
        que $(v'_j)_B$ es la columna $j$ de $P$. Las columnas de $P$ son
        una base de $k^n$ (pues $\rg P=n$, por el Tema~2). Por la
        traducción a coordenadas del Tema~2, $C=\{v'_1,\ldots,v'_n\}$ es
        una base de $V$. Por construcción, la columna $j$ de $M(C,B)$ es
        $(v'_j)_B$, así que $M(C,B)=P$.
    \end{proof}
\end{frame}


\begin{frame}
    \frametitle{Todas las matrices de un homomorfismo}

    \begin{theorem}
        Sea $f\colon V\to V'$ un homomorfismo y sea
        $A=M_{B,B'}(f)\in M_{m\times n}(k)$. Una matriz
        $A'\in M_{m\times n}(k)$ es la matriz de $f$ respecto de otras
        bases $C$ de $V$ y $C'$ de $V'$ si y solo si
        \[
            A'=QAP\quad\text{con }P\in M_n(k),\ Q\in M_m(k)\text{ invertibles.}
        \]
    \end{theorem}

    \begin{proof}
        ($\Rightarrow$) Es la fórmula del cambio de base, con
        $Q=M(B',C')$ y $P=M(C,B)$, que son invertibles (Tema~2).

        ($\Leftarrow$) Por el lema, existen bases $C$ de $V$ y $D$ de $V'$
        con $M(C,B)=P$ y $M(D,B')=Q^{-1}$. Sea $C'=D$; entonces
        $M(B',C')=M(D,B')^{-1}=Q$, y por la fórmula,
        \[
            M_{C,C'}(f)=M(B',C')\,A\,M(C,B)=QAP=A'.\qedhere
        \]
    \end{proof}
\end{frame}


\begin{frame}
    \frametitle{Matrices equivalentes y homomorfismos}

    Recordemos (Tema~1) que dos matrices $A,A'\in M_{m\times n}(k)$ son
    \alert{equivalentes por filas y columnas} si $A'=QAP$ con $P$ y $Q$
    invertibles, y que esto ocurre si y solo si $\rg A=\rg A'$.

    \begin{corollary}
        Sean $A,A'\in M_{m\times n}(k)$. Son equivalentes si y solo si
        existen un homomorfismo $f\colon V\to V'$ y bases $B,C$ de $V$ y
        $B',C'$ de $V'$ con
        \[
            A=M_{B,B'}(f),\qquad A'=M_{C,C'}(f).
        \]
    \end{corollary}

    \begin{proof}
        ($\Leftarrow$) Es el teorema anterior. ($\Rightarrow$) Si
        $A'=QAP$, tomamos $f=f_A\colon k^n\to k^m$ con las bases
        canónicas, de modo que $A=M_{B,B'}(f)$, y aplicamos el teorema.
    \end{proof}

    \medskip

    Dos matrices equivalentes son \alert{el mismo homomorfismo} visto en
    bases distintas, y la propiedad que comparten es su \alert{rango}.
\end{frame}


\begin{frame}
    \frametitle{Endomorfismos: una sola base}

    Si $f\colon V\to V$ es un endomorfismo, lo natural es usar la
    \alert{misma} base en el espacio de salida y en el de llegada.

    \begin{corollary}
        Sea $f\in\End(V)$, y sean $B$ y $C$ bases de $V$. Si
        $A=M_{B,B}(f)$, $A'=M_{C,C}(f)$ y $P=M(C,B)$, entonces
        \[
            A'=P^{-1}AP.
        \]
    \end{corollary}

    \begin{proof}
        Aplicamos la fórmula del cambio de base con $B'=B$ y $C'=C$, y
        usamos que $M(B,C)=M(C,B)^{-1}=P^{-1}$.
    \end{proof}

\end{frame}


\begin{frame}
    \frametitle{Matrices semejantes}

    \begin{definition}
        Dos matrices $A,A'\in M_n(k)$ son \alert{semejantes} si existe
        $P\in M_n(k)$ invertible con $A'=P^{-1}AP$.
    \end{definition}

    Las matrices de un mismo endomorfismo en bases distintas son
    semejantes (y, por el lema, recíprocamente). Encontrar una base en la
    que la matriz sea lo más sencilla posible es el problema central del
    Tema~5.
\end{frame}


\begin{frame}
    \frametitle{Completar una base del núcleo}

    \begin{lemma}
        Sea $f\colon V\to V'$ un homomorfismo con $V$ de dimensión finita.
        Sea $\{k_1,\ldots,k_m\}$ una base de $\ker f$ y sean
        $u_1,\ldots,u_r\in V$ tales que $f(u_1),\ldots,f(u_r)$ son una
        base de $\im f$. Entonces
        \[
            \{u_1,\ldots,u_r,k_1,\ldots,k_m\}
        \]
        es una base de $V$.
    \end{lemma}

    \begin{proof}
        Por el teorema de la dimensión, $r+m=\dim V$, así que basta ver que
        son linealmente independientes. Si
        $\sum_i\lambda_iu_i+\sum_j\mu_jk_j=0$, aplicando $f$ (y teniendo
        en cuenta que $f(k_j)=0$) obtenemos $\sum_i\lambda_if(u_i)=0$,
        luego todos los $\lambda_i$ son nulos. Entonces
        $\sum_j\mu_jk_j=0$ y todos los $\mu_j$ son nulos.
    \end{proof}

    \medskip

    En el Tema~2 vimos que $\dim V$ vectores independientes forman una
    base.
\end{frame}


\begin{frame}[allowframebreaks, t]
    \frametitle{Forma normal de un homomorfismo}

    \begin{theorem}
        Sea $f\colon V\to V'$ un homomorfismo entre espacios de dimensión
        finita, con $\rg f=r$. Existen bases $B$ de $V$ y $B'$ de $V'$ tales
        que
        \[
            M_{B,B'}(f)=\left(
            \begin{array}{c|c}
                I_r&0\\\hline0&0
            \end{array}
            \right).
        \]
    \end{theorem}

    \framebreak

    \begin{proof}
        Sea $\{k_1,\ldots,k_m\}$ una base de $\ker f$, con $m=n-r$.
        Como $\im f$ está generado por las imágenes de una base de $V$
        (Tema~2: podemos extraer una base de unos generadores), existen
        $u_1,\ldots,u_r\in V$ tales que $f(u_1),\ldots,f(u_r)$ son una base
        de $\im f$. Por el lema,
        \[
            B=\{u_1,\ldots,u_r,k_1,\ldots,k_m\}
        \]
        es una base de $V$. Ampliamos la familia independiente
        $f(u_1),\ldots,f(u_r)$ a una base de $V'$:
        \[
            B'=\{f(u_1),\ldots,f(u_r),w_{r+1},\ldots,w_{m'}\}.
        \]
        Entonces $f(u_i)$ es el $i$-ésimo vector de $B'$ para $i\leq r$, y
        $f(k_j)=0$. Por tanto, las primeras $r$ columnas de $M_{B,B'}(f)$ son
        $e_1,\ldots,e_r$ y las demás son nulas.
    \end{proof}
\end{frame}


\begin{frame}
    \frametitle{Forma normal y rango}

    \begin{corollary}
        Sea $f\colon V\to V'$ un homomorfismo entre espacios de dimensión
        finita.
        \begin{enumerate}[(1)]
            \item El rango de $M_{B,B'}(f)$ no depende de las bases elegidas
            y es igual a $\rg f$.
            \item Dos matrices $A,A'\in M_{m\times n}(k)$ son equivalentes
            si y solo si $\rg A=\rg A'$.
        \end{enumerate}
    \end{corollary}

    \begin{proof}
        (1) Ya lo sabíamos: $\rg M_{B,B'}(f)=\dim\im f$. (2) Si
        $\rg A=\rg A'=r$, ambas son la matriz de $f_A$ y de $f_{A'}$, que
        tienen rango $r$; por el teorema, cada una de ellas es equivalente
        a $\left(\begin{smallmatrix}I_r&0\\0&0\end{smallmatrix}\right)$.
        El recíproco es consecuencia de (1).
    \end{proof}

    \medskip

    Es el teorema de equivalencia por filas y columnas del Tema~1, ahora
    con una interpretación geométrica: \alert{existen bases en las que $f$
    es una proyección}, que lleva los primeros $r$ vectores de $B$ a los
    primeros $r$ de $B'$ y anula los demás.
\end{frame}


\begin{frame}
    \frametitle{Ejemplo: forma normal en $\R^4\to\R^3$}

    \begin{example}
        Para el homomorfismo $f\colon\R^4\to\R^3$ de antes, con $\rg f=2$,
        $\ker f=\langle(-2,1,0,0),(-1,0,-1,1)\rangle$ y
        $\im f=\langle(1,2,1),(0,1,1)\rangle$, busquemos bases con matriz
        $\left(\begin{smallmatrix}I_2&0\\0&0\end{smallmatrix}\right)$.
    \end{example}

    Los pivotes de $A$ estaban en las columnas $1$ y $3$, luego
    $f(e_1)=(1,2,1)$ y $f(e_3)=(0,1,1)$ son una base de $\im f$. Por el
    lema,
    \[
        B=\{e_1,\ e_3,\ (-2,1,0,0),\ (-1,0,-1,1)\}
    \]
    es una base de $\R^4$.
\end{frame}


\begin{frame}
    \frametitle{Ejemplo: forma normal en $\R^4\to\R^3$}

    Ampliamos $\{(1,2,1),(0,1,1)\}$ con $e_3$, que da una base de $\R^3$
    (la matriz de columnas es triangular con determinante $1$):
    \[
        B'=\{(1,2,1),\ (0,1,1),\ (0,0,1)\}.
    \]
    Como $f(e_1)$ y $f(e_3)$ son los dos primeros vectores de $B'$ y
    $f$ anula los dos últimos de $B$,
    \[
        M_{B,B'}(f)=\begin{pmatrix}1&0&0&0\\0&1&0&0\\0&0&0&0\end{pmatrix}.
    \]
\end{frame}


\begin{frame}
    \frametitle{Ejemplo: forma normal en polinomios}

    \begin{example}
        Para $\Delta\colon\R_2[x]\to\R_2[x]$, $\Delta(p)(x)=p(x+1)-p(x)$,
        con $\ker\Delta=\langle1\rangle$ y $\Delta(x)=1$,
        $\Delta(x^2)=1+2x$, busquemos bases con matriz normal.
    \end{example}

    Como $\Delta(x)=1$ y $\Delta(x^2)=1+2x$ son una base de $\im\Delta$,
    por el lema $\{x,x^2,1\}$ es una base de $\R_2[x]$. Ampliamos
    $\{1,1+2x\}$ con $x^2$ y tomamos
    \[
        B=\{x,\ x^2,\ 1\},\qquad B'=\{1,\ 1+2x,\ x^2\}.
    \]
    Entonces $\Delta(x)=1$ es el primer vector de $B'$, $\Delta(x^2)=1+2x$
    es el segundo y $\Delta(1)=0$, luego
    \[
        M_{B,B'}(\Delta)=\begin{pmatrix}1&0&0\\0&1&0\\0&0&0\end{pmatrix}.
    \]

    Con la \alert{misma} base de salida y de llegada esto es imposible
    para $\Delta$, pues $\Delta^3=0$ y $\operatorname{diag}(1,1,0)^3\neq0$.
\end{frame}
